\section{F9 · Partielle differensiallikninger og bølgelikningen}

\textbf{Kjernen:} Klassifiser PDE-en, still opp rand- og initialbetingelser, prøv $u=F(x)G(t)$ (separasjon av variable), finn egenfunksjonene fra randbetingelsene og bruk superposisjon og Fourierrekker for å tilpasse initialbetingelsene.

\subsection{Begreper og klassifisering}
\begin{defn}
En PDE er en likning med partiellderiverte av en ukjent funksjon $u$ av minst to variable, for eksempel $u(x,t)$ eller $u(x,y,t)$.
\begin{itemize}
  \item \textbf{Orden:} høyeste deriverte som forekommer.
  \item \textbf{Lineær:} $u$ og de deriverte forekommer bare i første potens og ikke multiplisert med hverandre. Koeffisientene kan avhenge av $x$ og $t$. Ellers er den \textbf{ikke-lineær}, for eksempel med $u^2$, $uu_x$ eller $\sin u$.
  \item \textbf{Homogen:} hvert ledd inneholder $u$ eller en derivert av $u$. Da er $u\equiv0$ en løsning. Et ledd uten $u$, som $f(x,t)$, en konstant $C\neq0$ eller $x$, gjør likningen \textbf{inhomogen}.
\end{itemize}
\end{defn}
\begin{eksempel}[øving 9, oppg.~1]
\begin{tabular}{@{}ll@{}}
$u_x+u^2=x$ & ikke-lineær ($u^2$), inhomogen ($x$), 1.\ orden\\
$u_t=u_{xx}$ & lineær, homogen, 2.\ orden (varmelikningen)\\
$\frac1r\partial_r(r\partial_ru)+\partial_\theta^2u=0$ & lineær, homogen (Laplace i polarkoordinater)\\
$u_{xx}+u_{yy}+C=0$ & lineær, inhomogen for $C\neq0$ (Poisson)\\
$u_t+u_x+xu^n=0$ & lineær hvis $n=1$, ellers ikke-lineær, homogen\\
$u_t+uu_x=u_{xx}$ & ikke-lineær ($uu_x$), homogen (Burgers)
\end{tabular}
\end{eksempel}

\begin{defn}[viktige likninger]
\begin{tabular}{@{}ll@{}}
Varmelikningen & $u_t-\alpha u_{xx}=q(x,t)$\\
Bølgelikningen & $u_{tt}=c^2u_{xx}$ ($c$ er bølgefarten), i 2D $u_{tt}=c^2(u_{xx}+u_{yy})$\\
Bjelke eller stav & $u_{tt}+k\,u_{xxxx}=q(x,t)$ (4.\ orden)\\
Laplace og Poisson & $\Delta u=u_{xx}+u_{yy}=0$ \ og \ $\Delta u=f$
\end{tabular}
\end{defn}

\begin{defn}[løsning, randbetingelser og initialbetingelser]
En løsning er en funksjon med de nødvendige deriverte som oppfyller likningen i området. For å få én bestemt løsning trengs \textbf{randbetingelser} (RB, verdier på randen av området) og \textbf{initialbetingelser} (IB, ved $t=0$). Med $\partial_t^k$ som høyeste tidsderiverte trengs $k$ initialbetingelser: én for varmelikningen og to (posisjon og hastighet) for bølgelikningen.
\end{defn}
\begin{teorem}[superposisjonsprinsippet]
Er $u_1$ og $u_2$ løsninger av en \emph{lineær homogen} PDE, så er $au_1+bu_2$ også en løsning. For en inhomogen likning gjelder ikke dette, men løsningen er en partikulær løsning pluss en homogen løsning (øving 9, oppg.~2).
\end{teorem}

\subsection{Bølgelikningen}
Modell: en streng på $[0,L]$ som er festet i endene. Strengen er homogen, tyngdekraften neglisjeres, og hvert punkt beveger seg bare loddrett. Bølgefarten er $c^2=\frac{\text{strekkraft}}{\text{masse per lengde}}$, så $c^2=\frac{FL}m$ i øving 9.

\begin{formel}[{Bølgelikningen på $[0,L]$ og løsningen\fs}]
\[
\begin{cases}
u_{tt}=c^2u_{xx}, & 0<x<L,\ t>0,\\
u(0,t)=u(L,t)=0 & \text{(RB)},\\
u(x,0)=f(x),\quad u_t(x,0)=g(x) & \text{(IB: posisjon og hastighet)},
\end{cases}
\]
\begin{gather*}
u(x,t)=\sum_{n=1}^\infty\Bigl(A_n\cos\frac{cn\pi t}{L}+B_n\sin\frac{cn\pi t}L\Bigr)\sin\frac{n\pi x}L,\\
A_n=\frac2L\int_0^Lf(x)\sin\frac{n\pi x}L\dd x,\qquad B_n=\frac2{cn\pi}\int_0^Lg(x)\sin\frac{n\pi x}L\dd x.
\end{gather*}
\end{formel}

\begin{oppskrift}[separasjon av variable (utledningen må kunnes)]
\begin{enumerate}
  \item \textbf{Ansats:} $u(x,t)=F(x)G(t)$. Innsatt gir det $FG''=c^2F''G$, altså
  \[
  \frac{F''(x)}{F(x)}=\frac{G''(t)}{c^2G(t)}=k .
  \]
  $k$ er en konstant, fordi venstre side bare avhenger av $x$ og høyre bare av $t$.
  Det gir to ODE-er: $F''-kF=0$ og $G''-c^2kG=0$. Randbetingelsene gir $F(0)=F(L)=0$ (ellers er $G\equiv0$).
  \item \textbf{Tilfellet $k=0$:} $F=Ax+B$, og RB gir $F\equiv0$.
  \item \textbf{Tilfellet $k=p^2>0$:} $F=A\ee^{px}+B\ee^{-px}$, og RB gir $A=B=0$.
  \item \textbf{Tilfellet $k=-p^2<0$:} $F=A\cos px+B\sin px$. $F(0)=0$ gir $A=0$, og $F(L)=0$ gir $\sin pL=0$, så $p=\frac{n\pi}L$. Altså er $F_n=\sin\frac{n\pi x}L$.
  \item \textbf{Tidsdelen:} $G_n''+\lambda_n^2G_n=0$ med $\lambda_n=\frac{cn\pi}L$, så $G_n=A_n\cos\lambda_nt+B_n\sin\lambda_nt$.
  \item \textbf{Superposisjon:} $u=\sum_nG_n(t)F_n(x)$. Sett $t=0$: $f(x)=\sum A_n\sin\frac{n\pi x}L$ er sinusrekken til $f$ (odde utvidelse). Deriver i $t$: $g(x)=\sum B_n\lambda_n\sin\frac{n\pi x}L$.
\end{enumerate}
\textbf{Begreper:} $u_n$ er egenfunksjonene, $\lambda_n=\frac{cn\pi}L$ egenverdiene, og $\{\lambda_n\}$ spekteret. Grunntonen har periode $\frac{2\pi}{\lambda_1}=\frac{2L}c$.
\end{oppskrift}

\begin{fallgruve}[effektive triks]
\begin{itemize}
  \item \textbf{Les av:} er $f$ en endelig sum av $\sin\frac{n\pi x}L$, er $A_n$ bare koeffisientene. For eksempel har $f=\sin(4\pi x)+3\sin(7\pi x)$ på $[0,1]$ $A_4=1$, $A_7=3$ og ellers $0$, så du slipper å integrere.
  \item \textbf{Er $g\equiv0$,} er alle $B_n=0$.
  \item \textbf{Et annet intervall,} for eksempel $[-L,L]$ med $u(\pm L,t)=0$ (prøveeksamen, oppg.~8): bruk $F=A\cos px+B\sin px$. Da gir $A\cos pL=0$ og $B\sin pL=0$ to familier: $\sin\frac{n\pi x}L$ og $\cos\frac{(2m-1)\pi x}{2L}$. Alternativt kan du forskyve til $[0,2L]$.
  \item \textbf{Neumann-rand} ($u_x=0$ i endene) gir $\cos\frac{n\pi x}L$, inkludert en konstant mode for $n=0$.
\end{itemize}
\end{fallgruve}

\begin{eksempel}[strengen løftes på midten (forelesningen)]
$f(x)=H\bigl(1-\abs{\frac{2x}L-1}\bigr)$ (trekant med topp $H$ i $x=\frac L2$) og $g=0$. Da er $B_n=0$ og $A_n=\frac{8H}{n^2\pi^2}\sin\frac{n\pi}{2}$, så bare odde $n$ bidrar, med vekslende fortegn.
\end{eksempel}

\begin{oving}
\begin{itemize}[leftmargin=5.6em,labelwidth=5em,labelsep=0.6em,align=left]
\item[\lett] \href{https://github.com/gustavbv/Matte4C---eksamen/blob/main/oppgavesett/F09-pde-og-bolgelikning/1-TMA4432-ovinger-og-proveeksamen/003-TMA4432-oving-9-oppg1-oppgave.pdf}{Øving 9, oppg.~1} {\color{graa}(\href{https://github.com/gustavbv/Matte4C---eksamen/blob/main/oppgavesett/F09-pde-og-bolgelikning/1-TMA4432-ovinger-og-proveeksamen/003-TMA4432-oving-9-oppg1-svar.pdf}{LF})}: klassifisering av PDE-er
\item[\middels] \href{https://github.com/gustavbv/Matte4C---eksamen/blob/main/oppgavesett/F09-pde-og-bolgelikning/2-Bergmann-eksamener/011-TMA4125-2024-var-oppg4-oppgave.pdf}{Matte 4N vår 2024, oppg.~4} {\color{graa}(\href{https://github.com/gustavbv/Matte4C---eksamen/blob/main/oppgavesett/F09-pde-og-bolgelikning/2-Bergmann-eksamener/011-TMA4125-2024-var-oppg4-svar.pdf}{LF})}: klassifisering og verifisering
\item[\middels] \href{https://github.com/gustavbv/Matte4C---eksamen/blob/main/oppgavesett/F09-pde-og-bolgelikning/1-TMA4432-ovinger-og-proveeksamen/001-TMA4432-2026-proveeksamen-oppg8-oppgave.pdf}{Prøveeksamen, oppg.~8} {\color{graa}(\href{https://github.com/gustavbv/Matte4C---eksamen/blob/main/oppgavesett/F09-pde-og-bolgelikning/1-TMA4432-ovinger-og-proveeksamen/001-TMA4432-2026-proveeksamen-oppg8-svar.pdf}{LF})}: separasjon på $[-1,1]$: sinus- og cosinusmoder
\item[\middels] \href{https://github.com/gustavbv/Matte4C---eksamen/blob/main/oppgavesett/F09-pde-og-bolgelikning/3-hele-oppgaven-i-pensum/023-TMA4130-35-2025-kont-oppg6-oppgave.pdf}{TMA4130 kont.\ 2025, oppg.~6} {\color{graa}(\href{https://github.com/gustavbv/Matte4C---eksamen/blob/main/oppgavesett/F09-pde-og-bolgelikning/3-hele-oppgaven-i-pensum/023-TMA4130-35-2025-kont-oppg6-svar.pdf}{LF})}: streng sluppet fra trekantform
\item[\vanskelig] \href{https://github.com/gustavbv/Matte4C---eksamen/blob/main/oppgavesett/F09-pde-og-bolgelikning/1-TMA4432-ovinger-og-proveeksamen/005-TMA4432-oving-9-oppg3-oppgave.pdf}{Øving 9, oppg.~3} {\color{graa}(\href{https://github.com/gustavbv/Matte4C---eksamen/blob/main/oppgavesett/F09-pde-og-bolgelikning/1-TMA4432-ovinger-og-proveeksamen/005-TMA4432-oving-9-oppg3-svar.pdf}{LF})}: trapesformet startprofil og svingeperiode
\item[\vanskelig] \href{https://github.com/gustavbv/Matte4C---eksamen/blob/main/oppgavesett/F09-pde-og-bolgelikning/2-Bergmann-eksamener/017-TMA4125-2022-var-oppg9-oppgave.pdf}{Matte 4N vår 2022, oppg.~9} {\color{graa}(\href{https://github.com/gustavbv/Matte4C---eksamen/blob/main/oppgavesett/F09-pde-og-bolgelikning/2-Bergmann-eksamener/017-TMA4125-2022-var-oppg9-svar.pdf}{LF})}: separasjon for en fjerdeordens PDE
\end{itemize}
\end{oving}

